\subsection{Continuous random variables}

A continuous random variable takes values over a continuous range, such as the real line or an interval on the real line.  
In the discrete case, one unit of probability is distributed over the possible values as point masses, and the PMF tells us how much probability sits on top of each value.  
In the continuous case, one unit of probability is still assigned to the possible values, but now it is spread continuously over the real line.

Some parts of the real line receive more probability mass per unit length, others receive less.  
The \textit{probability density function} (PDF), denoted \(f_X(x)\), describes how densely probability is distributed along the real line.  
The letter \(f\) indicates we are dealing with a PDF, and the subscript indicates which random variable we are considering.

\bigskip\textsc{Interpretation of the PDF}

Think of probability as snowfall.  
One pound of snow has fallen on the real line, and the PDF tells us the height of the snow accumulated over each point.  
The weight of snow sitting on top of an interval \([a, b]\) is found by calculating the area under the PDF curve over that interval.  
Mathematically, area under the curve is an integral, so
\[
\mathbb{P}(a \le X \le b) = \int_a^b f_X(x) \, dx.
\]
This is the probability that \(X\) takes values in the interval \([a, b]\).

\bigskip\textsc{Properties of a PDF}

By analogy with the discrete case, a PDF must satisfy two conditions:
\[
f_X(x) \ge 0 \quad \text{for all } x,
\]
\[
\int_{-\infty}^{+\infty} f_X(x) \, dx = 1.
\]
The first condition ensures we do not obtain negative probabilities.  
The second condition ensures that the total probability over the entire real line equals one, since \(X\) is certain to lie somewhere on the real line.  
These two conditions are all that we need for a legitimate PDF.

\bigskip\textsc{Formal definition of a continuous random variable}

A continuous random variable is a random variable whose probabilities can be described by a PDF.  
That is, there exists a function \(f_X\) such that
\[
\mathbb{P}(a \le X \le b) = \int_a^b f_X(x) \, dx
\]
for any interval \([a, b]\).  
It is not enough for a random variable to take values on a continuous set.  
A continuous random variable must also be describable by a PDF according to this formula.

\bigskip\textsc{Probabilities of more complicated sets}

Once we have the probabilities of intervals, we can use additivity to calculate the probabilities of more complicated sets.  
For example, suppose we want the probability that \(X\) lies between \(1\) and \(3\) or between \(4\) and \(5\).  
This is the probability of a region consisting of two disjoint intervals.  
By additivity,
\[
\mathbb{P}(1 \le X \le 3 \text{ or } 4 \le X \le 5) = \mathbb{P}(1 \le X \le 3) + \mathbb{P}(4 \le X \le 5).
\]
Each of the probabilities on the right is calculated using the PDF:
\[
= \int_1^3 f_X(x) \, dx + \int_4^5 f_X(x) \, dx.
\]

\bigskip\textsc{The role of the sample space}

There is still an underlying sample space lurking in the background.  
Different outcomes in the sample space result in different numerical values for the random variable \(X\).  
When we talk about the probability that \(X\) takes values in an interval \([a, b]\), we mean the probability of those outcomes for which the resulting value of \(X\) lies in that interval.

However, once we have a PDF in our hands, we can completely forget about the sample space and carry out any calculations just by working with the PDF.  
As we move on in this course, the sample space will appear less and less, and we will work directly with PDFs or PMFs.

\newpage
\bigskip\textsc{Probabilities of small intervals}

Let us look at an interval that starts at some point \(a\) and goes up to \(a + \delta\), where \(\delta\) is a small positive number.  
The probability that \(X\) falls in this interval is the area under the PDF over the interval \([a, a + \delta]\).  
As long as \(f_X\) does not change too much over this little interval, which will be the case if \(f_X\) is continuous, we can approximate the area by the area of a rectangle with constant height.  
The area of this rectangle is
\[
\mathbb{P}(a \le X \le a + \delta) \approx f_X(a) \cdot \delta.
\]
This gives us an interpretation of PDFs in terms of probabilities of small intervals.

If we rearrange this approximation, we see that
\[
f_X(a) \approx \frac{\mathbb{P}(a \le X \le a + \delta)}{\delta}.
\]
So the value of the PDF can be interpreted as \textit{probability per unit length}.  
PDFs are not probabilities, they are densities, and their units are probability per unit length.

\bigskip\textsc{Probability of a single point}

If the probability per unit length is finite and the length \(\delta\) is sent to zero, we obtain zero probability.  
More formally, if we let \(b = a\) in the integral
\[
\mathbb{P}(a \le X \le a) = \int_a^a f_X(x) \, dx,
\]
we get an integral over a zero-length interval, and that integral is zero.  
Therefore,
\[
\mathbb{P}(X = a) = 0
\]
for any particular point \(a\).

So for a continuous random variable, any particular point has zero probability.  
Yet, collectively, the infinitely many points in an interval together have positive probability.  
This is not a puzzle, it is exactly what happens with the ordinary notion of length.  
Single points have zero length, yet by putting together lots of points we can create a set with positive length.

\bigskip\textsc{Open and closed intervals}

A final consequence of the fact that individual points have zero probability is that it does not matter whether we include the endpoints of an interval or not.  
Using additivity,
\[
\mathbb{P}(a \le X \le b) = \mathbb{P}(X = a) + \mathbb{P}(a < X < b) + \mathbb{P}(X = b).
\]
Since \(\mathbb{P}(X = a) = 0\) and \(\mathbb{P}(X = b) = 0\), we conclude that
\[
\mathbb{P}(a \le X \le b) = \mathbb{P}(a < X < b).
\]
When calculating probabilities of intervals, it does not matter whether we include the endpoints or not.


\subsubsection{Uniform and piecewise constant PDFs}

The uniform random variable is a basic example of a continuous random variable.  
In the discrete uniform case, the possible values are integers between \(a\) and \(b\), all equally likely.  
In the continuous case, any real number between \(a\) and \(b\) is possible, and the PDF has the same height at all points in this interval.

Intuitively, a uniform random variable models a situation where we know that the value lies between \(a\) and \(b\), but we have no reason to believe that any location inside the interval is more likely than another.  
In this sense, it represents a situation of complete ignorance.

\bigskip\textsc{PDF of a uniform random variable}

Let \(X\) be uniform on \([a,b]\).  
Its PDF is constant on \([a,b]\) and zero elsewhere:
\[
f_X(x) =
\begin{cases}
c, & a \le x \le b, \\
0, & \text{otherwise}.
\end{cases}
\]
The total probability must be \(1\), so the area under the PDF is
\[
1 = \int_a^b c \, dx = c(b-a),
\]
which gives \(c = \frac{1}{b-a}\).  
Thus
\[
f_X(x) =
\begin{cases}
\frac{1}{b-a}, & a \le x \le b, \\
0, & \text{otherwise}.
\end{cases}
\]
Any two intervals of the same length inside \([a,b]\) have the same probability, because their probabilities are the base (length) times the same constant height.

\bigskip\textsc{Piecewise constant PDFs}

A more general kind of PDF is piecewise constant.  
In this case, the PDF is constant on each of several subintervals, possibly with different heights on different intervals.  
This shows that PDFs need not be continuous; they may have jumps and discontinuities.

For such a PDF to be valid, the total area under the curve must be \(1\).  
This total area is the sum of the areas of the rectangles corresponding to each constant piece.

\bigskip\textsc{Computing probabilities}

With a piecewise constant PDF, probabilities are obtained by adding areas of rectangles.  
If a certain event corresponds to an interval that intersects several constant pieces, its probability is the sum of the areas over each part of that interval.  
For example, if over two subintervals the heights are \(h_1\) and \(h_2\) and the lengths are \(\ell_1\) and \(\ell_2\), then the probability of their union is
\[
h_1 \ell_1 + h_2 \ell_2.
\]


\subsubsection{Means and variances}

We now develop continuous counterparts of the concepts of expectation and variance.  
In the discrete case, expectation is a weighted average of the possible values of a random variable, weighted by their probabilities.  
In the continuous case, we use the same idea, but the sum is replaced by an integral and the weights come from the PDF.

\bigskip\textsc{Expectation of a continuous random variable}

Let \(X\) be a continuous random variable with PDF \(f_X(x)\).  
Its expectation (or mean) is defined by
\[
\mathbb{E}[X] = \int_{-\infty}^{+\infty} x f_X(x) \, dx.
\]
This is a weighted average of the possible values of \(X\), where points with higher density \(f_X(x)\) contribute more.  
More generally, for any function \(g\),
\[
\mathbb{E}[g(X)] = \int_{-\infty}^{+\infty} g(x) f_X(x) \, dx.
\]
This is the continuous analogue of the discrete expected value rule.  

To ensure that the expectation is well defined, we usually assume that
\[
\int_{-\infty}^{+\infty} |x| f_X(x) \, dx < \infty.
\]
Under this assumption, \(\mathbb{E}[X]\) exists.  

The intuition of expectation remains the same as in the discrete case.  
It represents the average value we would see if we repeated the experiment many times independently.  
It can also be viewed as the center of gravity of the distribution.  
If the PDF is symmetric around a point \(c\), then the expectation is \(c\).

\bigskip\textsc{Basic properties of expectation}

Expectations of continuous random variables satisfy the same properties as in the discrete case.  
If \(X \ge 0\) always, then \(\mathbb{E}[X] \ge 0\).  
If \(X\) always lies in an interval \([a,b]\), then \(\mathbb{E}[X]\) also lies in \([a,b]\).

A key property is linearity of expectation.  
For constants \(a\) and \(b\),
\[
\mathbb{E}[aX + b] = a\,\mathbb{E}[X] + b.
\]
Linearity holds regardless of whether \(X\) is discrete or continuous.  

\bigskip\textsc{Variance and standard deviation}

The variance of a continuous random variable is defined exactly as in the discrete case:
\[
\mathrm{Var}(X) = \mathbb{E}\big[(X - \mu)^2\big],
\]
where \(\mu = \mathbb{E}[X]\).  
Using the expected value rule, this becomes
\[
\mathrm{Var}(X) = \int_{-\infty}^{+\infty} (x - \mu)^2 f_X(x) \, dx.
\]
The standard deviation is
\[
\sigma = \sqrt{\mathrm{Var}(X)}.
\]

Variance measures how spread out the values of \(X\) are around the mean.  
A small variance means that most of the probability mass is concentrated close to \(\mu\).  
A large variance means that values far from \(\mu\) occur more often.

\bigskip\textsc{Alternative formula for the variance}

It is often convenient to use the formula
\[
\mathrm{Var}(X) = \mathbb{E}[X^2] - \big(\mathbb{E}[X]\big)^2.
\]
Here
\[
\mathbb{E}[X^2] = \int_{-\infty}^{+\infty} x^2 f_X(x) \, dx.
\]
This alternative expression is derived exactly as in the discrete case and is usually easier to use in calculations.

\bigskip\textsc{Effect of linear transformations}

If we transform a random variable \(X\) by a linear function \(aX + b\), the variance behaves as follows:
\[
\mathrm{Var}(aX + b) = a^2 \mathrm{Var}(X).
\]
Adding a constant \(b\) does not change the variance.  
Multiplying by a constant \(a\) scales the variance by \(a^2\).


\subsubsection{Mean and variance of the uniform}

We now illustrate mean and variance calculations using the continuous uniform random variable.  
This is the continuous analogue of the discrete uniform, for which we already know formulas for the mean and variance.

\bigskip\textsc{Mean of the uniform}

Let \(X\) be uniform on \([a,b]\), with density
\[
f_X(x) =
\begin{cases}
\frac{1}{b-a}, & a \le x \le b, \\
0, & \text{otherwise}.
\end{cases}
\]
The mean is
\[
\mathbb{E}[X] = \int_{-\infty}^{+\infty} x f_X(x)\,dx.
\]
Since \(f_X(x)=0\) outside \([a,b]\), this reduces to
\[
\mathbb{E}[X] = \int_a^b x \cdot \frac{1}{b-a}\,dx = \frac{1}{b-a}\int_a^b x\,dx.
\]
We have
\[
\int_a^b x\,dx = \frac{b^2}{2} - \frac{a^2}{2},
\]
so
\[
\mathbb{E}[X] = \frac{1}{b-a}\left(\frac{b^2}{2} - \frac{a^2}{2}\right) = \frac{a+b}{2}.
\]
This is the same as in the discrete case and also matches the midpoint of the interval \([a,b]\), reflecting the symmetry of the PDF around \(\frac{a+b}{2}\).

\bigskip\textsc{Computing \(\mathbb{E}[X^2]\)}

To use the alternative variance formula, we first compute \(\mathbb{E}[X^2]\).  
By the expected value rule,
\[
\mathbb{E}[X^2] = \int_{-\infty}^{+\infty} x^2 f_X(x)\,dx = \int_a^b x^2 \cdot \frac{1}{b-a}\,dx.
\]
Thus
\[
\mathbb{E}[X^2] = \frac{1}{b-a}\int_a^b x^2\,dx.
\]
Since
\[
\int_a^b x^2\,dx = \frac{b^3}{3} - \frac{a^3}{3},
\]
we obtain
\[
\mathbb{E}[X^2] = \frac{1}{b-a}\left(\frac{b^3}{3} - \frac{a^3}{3}\right).
\]
The cubic terms appear because the antiderivative of \(x^2\) is \(\frac{x^3}{3}\).

\bigskip\textsc{Variance and standard deviation}

Using the alternative formula
\[
\mathrm{Var}(X) = \mathbb{E}[X^2] - \big(\mathbb{E}[X]\big)^2,
\]
we substitute
\[
\mathbb{E}[X^2] = \frac{1}{b-a}\left(\frac{b^3}{3} - \frac{a^3}{3}\right),
\quad
\mathbb{E}[X] = \frac{a+b}{2}.
\]
After some algebra, this simplifies to
\[
\mathrm{Var}(X) = \frac{(b-a)^2}{12}.
\]
The standard deviation is then
\[
\sigma = \sqrt{\mathrm{Var}(X)} = \frac{b-a}{\sqrt{12}}.
\]

A useful observation is that the standard deviation is proportional to the width \(b-a\) of the interval.  
The wider the uniform distribution, the larger the standard deviation, which matches the intuition that standard deviation measures the spread or width of the distribution.  
The variance grows with the square of the length \(b-a\), reflecting that doubling the width multiplies the variance by four.


\subsubsection{Exponential random variables}

We now introduce a new continuous random variable, the exponential random variable.  
It is determined by a single positive parameter \(\lambda > 0\), and its PDF is
\[
f_X(x) =
\begin{cases}
\lambda e^{-\lambda x}, & x \ge 0, \\
0, & x < 0.
\end{cases}
\]
Since the PDF is zero for negative \(x\), the exponential random variable is always non-negative.

For \(x \ge 0\), the PDF starts at value \(\lambda\) when \(x = 0\) and then decays exponentially.  
When \(\lambda\) is small, the initial height is small but the decay is slow, so the distribution extends over a wide range.  
When \(\lambda\) is large, the PDF starts high near zero and decays quickly, so most of the probability is concentrated near \(0\).

The shape of the exponential PDF is analogous to that of the geometric distribution in the discrete case.  
There, probabilities decrease geometrically with the value; here they decrease exponentially.

\bigskip\textsc{Tail probabilities}

Fix \(a > 0\).  
We want to compute
\[
\mathbb{P}(X \ge a) = \int_a^{\infty} \lambda e^{-\lambda x} \, dx.
\]
Using the antiderivative \(\int e^{-\lambda x} dx = -\frac{1}{\lambda} e^{-\lambda x}\), we get
\[
\mathbb{P}(X \ge a)
= \lambda \left[-\frac{1}{\lambda} e^{-\lambda x}\right]_{x=a}^{x=\infty}
= - e^{-\lambda x}\Big|_{x=a}^{x=\infty}
= 0 - \big(-e^{-\lambda a}\big)
= e^{-\lambda a}.
\]
Thus the tail probability decays exponentially in \(a\).  

If we set \(a = 0\), we obtain
\[
\mathbb{P}(X \ge 0) = e^{-\lambda \cdot 0} = 1,
\]
which confirms that the total integral of the PDF over its support is \(1\).

\bigskip\textsc{Mean and second moment}

To compute the mean, we use the definition
\[
\mathbb{E}[X] = \int_0^{\infty} x \lambda e^{-\lambda x} \, dx.
\]
This integral can be evaluated by integration by parts and the result is
\[
\mathbb{E}[X] = \frac{1}{\lambda}.
\]

For the second moment, we similarly write
\[
\mathbb{E}[X^2] = \int_0^{\infty} x^2 \lambda e^{-\lambda x} \, dx,
\]
which is again computed by integration by parts (a bit more tedious).  
The result is
\[
\mathbb{E}[X^2] = \frac{2}{\lambda^2}.
\]

\bigskip\textsc{Variance}

Using the alternative variance formula,
\[
\mathrm{Var}(X) = \mathbb{E}[X^2] - \big(\mathbb{E}[X]\big)^2,
\]
we substitute
\[
\mathbb{E}[X^2] = \frac{2}{\lambda^2},
\quad
\mathbb{E}[X] = \frac{1}{\lambda}.
\]
Thus
\[
\mathrm{Var}(X)
= \frac{2}{\lambda^2} - \left(\frac{1}{\lambda}\right)^2
= \frac{2}{\lambda^2} - \frac{1}{\lambda^2}
= \frac{1}{\lambda^2}.
\]

When \(\lambda\) is small, the PDF decays slowly, so large values of \(X\) are more likely.  
In that case, both the mean \(\frac{1}{\lambda}\) and the variance \(\frac{1}{\lambda^2}\) are large, reflecting a wide and spread-out distribution.  

\bigskip\textsc{Interpretation and relation to geometric}

The exponential distribution is closely related to the geometric distribution.  
The mean \(\mathbb{E}[X] = 1/\lambda\) parallels the geometric mean \(1/p\), where \(p\) is the success probability.  
The exponential is often used to model waiting times until an event occurs, such as the time until a customer arrives, a machine fails, a light bulb burns out, or an email arrives.


\subsubsection{Cumulative distribution functions}

We now introduce a way to describe the distribution of any random variable, discrete or continuous, using the \textit{cumulative distribution function} (CDF).  
For a random variable \(X\), the CDF is the function
\[
F_X(x) = \mathbb{P}(X \le x).
\]
It gives the probability that \(X\) takes a value less than or equal to \(x\).

The CDF is defined the same way for all types of random variables.  
If \(X\) is continuous with PDF \(f_X\), then
\[
F_X(x) = \int_{-\infty}^{x} f_X(t)\,dt.
\]

\bigskip\textsc{Example: uniform random variable}

Let \(X\) be uniform on \([a,b]\).  
The CDF \(F_X(x)\) has three cases:
\[
F_X(x) =
\begin{cases}
0, & x \le a, \\
\displaystyle \frac{x-a}{b-a}, & a < x < b, \\
1, & x \ge b.
\end{cases}
\]
For \(x \le a\), there is no probability mass to the left, so the CDF is \(0\).  
For \(a < x < b\), the probability up to \(x\) is the area of a rectangle with base \(x-a\) and height \(\frac{1}{b-a}\), giving \(\frac{x-a}{b-a}\).  
For \(x \ge b\), we have already covered all the probability, so the CDF is \(1\).

The graph starts at \(0\), rises linearly from \(a\) to \(b\), and then stays at \(1\).

\bigskip\textsc{Using the CDF to get probabilities}

Once we know the CDF, we can compute probabilities of intervals.  
For example,
\[
\mathbb{P}(3 < X \le 4) = F_X(4) - F_X(3).
\]
More generally, for any \(a < b\),
\[
\mathbb{P}(a < X \le b) = F_X(b) - F_X(a).
\]
Thus the CDF contains all probabilistic information about the distribution.

For continuous random variables, we can recover the PDF from the CDF wherever the derivative exists:
\[
f_X(x) = \frac{d}{dx} F_X(x).
\]
Conversely, given the PDF, we obtain the CDF by integrating.

\bigskip\textsc{CDFs for discrete random variables}

For a discrete random variable \(X\), the CDF is still
\[
F_X(x) = \mathbb{P}(X \le x),
\]
but it is computed as a sum of PMF values for all outcomes \(k\) with \(k \le x\).  
The CDF has the shape of a staircase: it is constant between the possible values of \(X\) and jumps upward at each value where \(X\) has positive probability.

At a point \(x_0\) where \(X\) can take the value \(x_0\), the jump size is
\[
F_X(x_0) - \lim_{x \uparrow x_0} F_X(x) = \mathbb{P}(X = x_0).
\]
The value of the CDF at a jump is the value \emph{including} that point, because the definition uses \(\le x\).

\bigskip\textsc{General properties of CDFs}

CDFs satisfy some general properties for any random variable:
\[
\text{(i) } F_X(x) \text{ is non-decreasing in } x,
\]
\[
\text{(ii) } \lim_{x \to -\infty} F_X(x) = 0,
\]
\[
\text{(iii) } \lim_{x \to +\infty} F_X(x) = 1.
\]
The non-decreasing property reflects that larger \(x\) correspond to larger events \(\{X \le x\}\).  
As we move far to the left, the probability of being \(\le x\) must go to \(0\), and as we move far to the right, the probability must approach \(1\).

\subsubsection{Normal random variables}

Normal random variables, also called Gaussian random variables, are among the most important in probability theory.  
They play a key role in theoretical results such as the central limit theorem and are widely used in applications as models of random noise.

\bigskip\textsc{Standard normal distribution}

The \textit{standard normal} random variable is denoted by
\[
X \sim \mathcal{N}(0,1).
\]
Its PDF is defined for all real \(x\) by
\[
f_X(x) = \frac{1}{\sqrt{2\pi}} e^{-x^2/2}, \quad -\infty < x < \infty.
\]
This random variable can take any real value.  

The factor \(e^{-x^2/2}\) has a bell shape centered at \(0\): it is equal to \(1\) at \(x=0\) and decreases as \(|x|\) increases.  
The constant \(\frac{1}{\sqrt{2\pi}}\) is chosen so that the total area under the curve is \(1\).

Because the PDF is symmetric around \(0\), the mean of the standard normal is
\[
\mathbb{E}[X] = 0.
\]
A calculus computation shows that its variance is
\[
\mathrm{Var}(X) = 1.
\]

\bigskip\textsc{General normal distributions}

A general normal random variable with mean \(\mu\) and variance \(\sigma^2\) is denoted by
\[
X \sim \mathcal{N}(\mu,\sigma^2),
\]
where \(\sigma > 0\).  
Its PDF is
\[
f_X(x) = \frac{1}{\sqrt{2\pi}\,\sigma}
\exp\!\left(-\frac{(x-\mu)^2}{2\sigma^2}\right), \quad -\infty < x < \infty.
\]

The exponent \(-\frac{(x-\mu)^2}{2\sigma^2}\) is a quadratic centered at \(x=\mu\), so the PDF has a bell shape centered at \(\mu\).  
By symmetry, the mean is
\[
\mathbb{E}[X] = \mu.
\]
A calculus exercise shows that the variance is
\[
\mathrm{Var}(X) = \sigma^2.
\]

The parameter \(\sigma\) controls the spread.  
When \(\sigma\) is small, the quadratic in the exponent rises quickly, the bell is narrow, and the variance is small.  
When \(\sigma\) is large, the bell is wide and the variance is large.

\bigskip\textsc{Linear transformations of normal variables}

Normal random variables behave very nicely under linear transformations.  
Let
\[
X \sim \mathcal{N}(\mu,\sigma^2)
\]
and define
\[
Y = aX + b
\]
for constants \(a\) and \(b\).  

From linearity of expectation,
\[
\mathbb{E}[Y] = a\,\mathbb{E}[X] + b = a\mu + b.
\]
From the variance rule for linear functions,
\[
\mathrm{Var}(Y) = a^2 \mathrm{Var}(X) = a^2 \sigma^2.
\]
The key additional fact is that \(Y\) is itself normal:
\[
Y \sim \mathcal{N}(a\mu + b,\, a^2\sigma^2).
\]

There is one special case.  
If \(a = 0\), then \(Y = b\) is a constant (a degenerate discrete random variable with variance \(0\)).  
By convention, we can view this as a degenerate normal with mean \(b\) and variance \(0\), so that any linear function of a normal random variable is still considered normal.


\subsubsection{Calculation of normal probabilities}

Normal random variables are very important, so we would like to calculate probabilities such as \(\mathbb{P}(X \le 5)\) when \(X\) is normal.  
In general, the CDF of a normal random variable does not have a closed form expression in terms of elementary functions.  

\bigskip\textsc{Standard normal CDF and its table}

For the standard normal random variable \(Y \sim \mathcal{N}(0,1)\), we denote the CDF by
\[
\phi(y) = \mathbb{P}(Y \le y).
\]
Graphically, \(\phi(y)\) is the area under the standard normal PDF to the left of \(y\).  
Values of \(\phi(y)\) are tabulated in standard normal tables.

To use such a table, we locate the row and column corresponding to the first two decimal digits of \(y\).  
For example, \(\phi(0) = 0.5\) by symmetry of the standard normal around \(0\).  
If we want \(\phi(1.16)\), we find the row for \(1.1\) and the column for \(0.06\), and read off the corresponding entry.  
For \(y = 2.9\), we find the row for \(2.9\) and the column for \(0.00\), which gives \(\phi(2.9)\).  

Typical tables list \(\phi(y)\) only for \(y \ge 0\).  
To handle negative arguments, we use symmetry of the standard normal:
\[
\phi(-y) = \mathbb{P}(Y \le -y) = 1 - \mathbb{P}(Y \le y) = 1 - \phi(y), \quad y \ge 0.
\]
Thus, to find \(\phi(-2)\), we compute \(1 - \phi(2)\).  

\bigskip\textsc{Standardization of a general normal}

Let \(X\) be a normal random variable with mean \(\mu\) and standard deviation \(\sigma > 0\), that is \(X \sim \mathcal{N}(\mu,\sigma^2)\).  
We define the standardized random variable
\[
Y = \frac{X - \mu}{\sigma}.
\]
The quantity \(Y\) measures how many standard deviations \(X\) is away from its mean.  

The mean of \(Y\) is
\[
\mathbb{E}[Y] = \frac{\mathbb{E}[X] - \mu}{\sigma} = 0,
\]
and the variance of \(Y\) is
\[
\mathrm{Var}(Y) = \frac{1}{\sigma^2} \mathrm{Var}(X) = 1.
\]
If \(X\) is normal, then \(Y\) is a standard normal random variable.  
Equivalently,
\[
X = \mu + \sigma Y, \quad Y \sim \mathcal{N}(0,1).
\]

\bigskip\textsc{Using standardization to compute probabilities}

To compute probabilities for \(X\), we transform the event into one involving \(Y\).  
For example, for any real \(a\) and \(b\) with \(a < b\),
\[
\mathbb{P}(a \le X \le b)
= \mathbb{P}\!\left(a \le \mu + \sigma Y \le b\right)
= \mathbb{P}\!\left(\frac{a-\mu}{\sigma} \le Y \le \frac{b-\mu}{\sigma}\right).
\]
If we write
\[
\alpha = \frac{a-\mu}{\sigma}, \quad \beta = \frac{b-\mu}{\sigma},
\]
then
\[
\mathbb{P}(a \le X \le b) = \mathbb{P}(\alpha \le Y \le \beta) = \phi(\beta) - \phi(\alpha).
\]
We can then use the standard normal table to find \(\phi(\beta)\) and \(\phi(\alpha)\), handling negative arguments via symmetry if needed.

\bigskip\textsc{Example of the method}

Suppose \(X \sim \mathcal{N}(6,4)\), so \(\mu = 6\) and \(\sigma = 2\).  
We want
\[
\mathbb{P}(2 \le X \le 8).
\]
We standardize:
\[
\frac{2 - 6}{2} = -2, \quad \frac{8 - 6}{2} = 1.
\]
So
\[
\mathbb{P}(2 \le X \le 8) = \mathbb{P}(-2 \le Y \le 1),
\]
where \(Y \sim \mathcal{N}(0,1)\).  
This equals
\[
\mathbb{P}(Y \le 1) - \mathbb{P}(Y \le -2)
= \phi(1) - \phi(-2).
\]
Using symmetry,
\[
\phi(-2) = 1 - \phi(2),
\]
so
\[
\mathbb{P}(2 \le X \le 8)
= \phi(1) - \big(1 - \phi(2)\big)
= \phi(1) + \phi(2) - 1.
\]
Finally, we read \(\phi(1)\) and \(\phi(2)\) from the standard normal table and substitute the corresponding numerical values.